\documentclass{article}

% ICML 2026 style (use official icml2026.sty when available)
\usepackage[utf8]{inputenc}
\usepackage{amsmath,amsfonts,amssymb}
\usepackage{graphicx}
\usepackage{booktabs}
\usepackage{hyperref}
\usepackage{algorithm}
\usepackage{multirow}
\usepackage{algpseudocode}
\usepackage{xcolor}
\usepackage[margin=1in]{geometry}

\title{Compositional Verification of Robust Mixture-of-Experts Architectures}

\author{
Connor Pham \\
Vanderbilt University \\
\texttt{connor.pham@vanderbilt.edu}
}

\date{}

\begin{document}

\maketitle

%==============================================================================
% ABSTRACT 
%==============================================================================
\begin{abstract}
Mixture-of-Experts (MoE) architectures offer modularity and scalability, yet their formal verifiability remains poorly understood. We present a \textbf{compositional verification framework} for heterogeneous MoE systems: if the router provably maintains consistent expert selection within an adversarial perturbation region, and the selected expert is provably robust, then the entire MoE system is provably robust. This decomposition reduces verification complexity from the full system to independent component verification.

Empirically, we demonstrate that (1) expert robustness dominates system behavior---adversarially trained experts improve system adversarial accuracy by 13.1 percentage points while router training shows negligible impact ($<$0.1\%); (2) verified routers achieve 100\% certified accuracy at practical threat models ($\varepsilon \leq 4/255$); and (3) the compositional property enables 20$\times$ parameter reduction for verification. These results establish a principled path toward scalable, formally verified modular AI systems.
\end{abstract}

%==============================================================================
% 1. INTRODUCTION 
%==============================================================================
\section{Introduction}

Deep neural networks achieve state-of-the-art performance but remain vulnerable to adversarial perturbations---a critical liability for safety-critical applications. Mixture-of-Experts (MoE) architectures have emerged as a powerful paradigm for scaling models through modular computation, yet their robustness properties, particularly in heterogeneous settings where domain-specific experts cooperate, are not well understood.

Formal verification tools like $\alpha$--$\beta$-CROWN now enable provable robustness guarantees for neural networks. However, verifying full MoE systems is intractable due to their scale, sparsity, and dynamic routing. A key question is whether robustness \emph{composes}: can component-level verification yield system-level guarantees?

We answer affirmatively by introducing \textbf{MetaMoE}, a verification-aware heterogeneous MoE framework. Our central contribution is a \textbf{compositional verification theorem}:

\begin{quote}
\textit{If the router maintains invariant expert selection within an $\ell_\infty$ perturbation ball, and the selected expert is robust within that ball, then the MoE system is robust.}
\end{quote}

This decomposition enables modular verification: verify the router and experts independently, then compose guarantees without re-verifying the full system. Our contributions are:

\begin{enumerate}
    \item \textbf{Compositional Verification Framework:} We formalize and validate conditions under which MoE robustness decomposes into router stability and expert robustness, enabling scalable modular certification.
    
    \item \textbf{Expert Dominance Finding:} We demonstrate that expert training determines system robustness (13.1\% adversarial accuracy difference), while router adversarial training has negligible impact ($<$0.1\%) for visually dissimilar domains---enabling 4$\times$ training speedup.
    
    \item \textbf{Empirical-Formal Correlation:} We quantify domain-dependent certification gaps and show robust training simultaneously improves empirical and certified accuracy.
    
    \item \textbf{Verification Efficiency:} Compositional verification achieves 20$\times$ parameter reduction (2M $\to$ 96K) and enables incremental expert integration without full-system re-verification.
\end{enumerate}

%==============================================================================
% 2. BACKGROUND AND RELATED WORK 
%==============================================================================
\section{Background and Related Work}

\paragraph{Mixture-of-Experts.} The MoE paradigm \cite{jacobs1991} decomposes problems into specialized subtasks coordinated by gating networks. Shazeer et al.~\cite{shazeer2017} scaled this via sparse top-$k$ routing, and Riquelme et al.~\cite{riquelme2021} extended MoE to vision. Recent work explores heterogeneous compositions \cite{fedus2023} and modular architectures \cite{pfeiffer2023}, but robustness properties remain underexplored.

\paragraph{Adversarial Robustness.} Adversarial training \cite{madry2018} remains the primary defense against $\ell_p$-bounded attacks. The robustness-accuracy tradeoff is well-documented, with robust models sacrificing clean accuracy for adversarial resilience.

\paragraph{Neural Network Verification.} Complete verifiers like $\alpha$--$\beta$-CROWN \cite{wang2021} provide sound robustness certificates via bound propagation and branch-and-bound. While effective for small networks, scalability to large systems remains challenging. Compositional verification---verifying components independently and composing guarantees---offers a path forward but requires careful analysis of component interactions.

\paragraph{Research Gap.} Prior work focuses on homogeneous MoE scalability or monolithic verification. We address the intersection: \textit{compositional robustness verification for heterogeneous MoE systems}.

%==============================================================================
% 3. COMPOSITIONAL VERIFICATION FRAMEWORK
%==============================================================================
\section{Compositional Verification Framework}

\subsection{Problem Setup}

Let $x \in \mathbb{R}^{C \times H \times W}$ be an input with label $y$. A heterogeneous MoE consists of:
\begin{itemize}
    \item \textbf{Experts:} $\{E_i: \mathbb{R}^d \to \mathbb{R}^{C_i}\}_{i=1}^N$ producing domain-specific logits
    \item \textbf{Router:} $\mathcal{G}: \mathbb{R}^d \to \mathbb{R}^N$ selecting expert $m^* = \arg\max_i \mathcal{G}_i(x)$
\end{itemize}

The MoE output combines selected expert predictions in a unified class space via offsets $o_i = \sum_{j<i} C_j$.

\paragraph{Threat Model.} We consider $\ell_\infty$-bounded perturbations: $x' \in B_\varepsilon(x) = \{x': \|x' - x\|_\infty \leq \varepsilon\}$.

\paragraph{Robustness Definitions.}
\begin{itemize}
    \item \textbf{Expert Robustness:} $E_i$ is $\varepsilon$-robust at $x$ if $\arg\max E_i(x') = y$ for all $x' \in B_\varepsilon(x)$
    \item \textbf{Router Stability:} $\mathcal{G}$ is $\varepsilon$-stable at $x$ if $\arg\max_i \mathcal{G}_i(x') = m^*(x)$ for all $x' \in B_\varepsilon(x)$
    \item \textbf{MoE Robustness:} The MoE is $\varepsilon$-robust if the final prediction is invariant within $B_\varepsilon(x)$
\end{itemize}

\subsection{Compositional Robustness Theorem}

\begin{theorem}[Compositional Robustness]
\label{thm:compositional}
Let $\text{MoE} = (\mathcal{G}, \{E_i\})$ be a top-1 routing MoE. For input $x$ with router selection $m^* = \arg\max_i \mathcal{G}_i(x)$:
\begin{equation}
\underbrace{\mathcal{G} \text{ is } \varepsilon\text{-stable at } x}_{\text{Router verification}} \land \underbrace{E_{m^*} \text{ is } \varepsilon\text{-robust at } x}_{\text{Expert verification}} \implies \underbrace{\text{MoE is } \varepsilon\text{-robust at } x}_{\text{System guarantee}}
\end{equation}
\end{theorem}

\begin{proof}[Proof Sketch]
For any $x' \in B_\varepsilon(x)$: (1) router stability ensures $\arg\max_i \mathcal{G}_i(x') = m^*$, so expert $E_{m^*}$ is selected; (2) expert robustness ensures $\arg\max E_{m^*}(x') = y$. Since the MoE output depends only on the selected expert's prediction, the system prediction remains $y$.
\end{proof}

\paragraph{Verification Decomposition.} Theorem~\ref{thm:compositional} enables modular verification:
\begin{enumerate}
    \item Verify router: For each test input, check $\mathcal{G}_{m^*}(x') > \mathcal{G}_j(x')$ for all $j \neq m^*$, $x' \in B_\varepsilon(x)$
    \item Verify experts: For each expert $E_i$, verify robustness on its domain
    \item Compose: Verified router $\land$ verified expert $\implies$ verified MoE
\end{enumerate}

This reduces verification from the full MoE ($\sim$2M parameters in our setup) to the largest component ($\sim$96K parameters)---a 20$\times$ reduction.

\subsection{Verification-Aware Architecture}

Both router and experts use \textbf{VerifiableCNN} (96K parameters) designed for $\alpha$--$\beta$-CROWN compatibility: average pooling (linear, unlike max pooling), no BatchNorm (eliminated via weight folding), and minimal ONNX operators. Architecture: 4 conv blocks ($3 \to 20 \to 28 \to 40 \to 56$ channels), global average pooling, 2-layer FC head.

\paragraph{Training.} Non-Robust Training (NRT) uses standard cross-entropy. Robust Training (RT) employs PGD adversarial training ($\varepsilon = 8/255$, 7 steps, step size $2/255$).

%==============================================================================
% 4. EXPERIMENTS 
%==============================================================================
\section{Experiments}

We evaluate on CIFAR-10 (natural images) and MNIST (handwritten digits)---visually dissimilar domains enabling clean isolation of compositional effects. Extended experiments on similar domains (traffic sign datasets) appear in Appendix~\ref{app:similar_domains}.

\subsection{Router Formal Verification}

We verify router stability using $\alpha$--$\beta$-CROWN on 40 samples (20 per dataset) across $\varepsilon \in \{2/255, 4/255, 8/255\}$, averaged over 5 independent training runs.

\begin{table}[h]
\centering
\small
\begin{tabular}{llccc}
\toprule
\textbf{Training} & $\boldsymbol{\varepsilon}$ & \textbf{Verified} & \textbf{Falsified} & \textbf{Success Rate} \\
\midrule
\multirow{3}{*}{NRT} & 2/255 & 40/40 & 0 & 100\% \\
& 4/255 & 40/40 & 0 & 100\% \\
& 8/255 & 24/40 & 0 & 60\% \\
\midrule
\multirow{3}{*}{RT} & 2/255 & 40/40 & 0 & 100\% \\
& 4/255 & 40/40 & 0 & 100\% \\
& 8/255 & 40/40 & 0 & 100\% \\
\bottomrule
\end{tabular}
\caption{Router verification results. Both NRT and RT achieve 100\% certified accuracy at practical threat models ($\varepsilon \leq 4/255$). RT extends verification to $\varepsilon = 8/255$. No falsifications (property violations) in any configuration---all failures are timeouts.}
\label{tab:router_verification}
\end{table}

\textbf{Key Finding:} Routers achieve high formal certification even without adversarial training at practical threat models. The visual dissimilarity between CIFAR-10 and MNIST creates large decision margins ($\gg \varepsilon$), enabling robust routing. RT extends certified robustness to larger perturbations.

\subsection{Expert Formal Verification}

\begin{table}[h]
\centering
\small
\begin{tabular}{llcccc}
\toprule
\textbf{Expert} & \textbf{Training} & \multicolumn{3}{c}{\textbf{CRA at $\varepsilon$}} \\
& & 2/255 & 4/255 & 8/255 \\
\midrule
\multirow{2}{*}{CIFAR-10} & NRT & 0\% & 0\% & 0\% \\
& RT & 90\% & 54\% & 0\% \\
\midrule
\multirow{2}{*}{MNIST} & NRT & 75\% & 0\% & 0\% \\
& RT & 100\% & 99\% & 95\% \\
\bottomrule
\end{tabular}
\caption{Expert Certified Robustness Accuracy (CRA). RT enables verification: CIFAR-10-NRT is completely unverifiable, while CIFAR-10-RT achieves 90\% CRA at $\varepsilon=2/255$. MNIST-RT achieves near-complete certification.}
\label{tab:expert_verification}
\end{table}

\textbf{Key Finding:} Domain complexity drives verifiability---MNIST's simple features enable tight bounds, while CIFAR-10 requires robust training for any verification success. Robust training is a \textit{prerequisite} for verifiability, not merely helpful.

\subsection{Compositional MoE Robustness}

We evaluate all combinations of router (RT/NRT) and expert (RT/NRT for each domain) training, with 5 runs per configuration.

\begin{table}[h]
\centering
\small
\begin{tabular}{lcc}
\toprule
\textbf{Configuration} & \textbf{Clean Acc (\%)} & \textbf{Adv Acc (\%)} \\
\midrule
Both Experts RT & $87.98 \pm 0.06$ & $\mathbf{41.93 \pm 0.41}$ \\
Both Experts NRT & $\mathbf{91.32 \pm 0.06}$ & $28.78 \pm 0.59$ \\
CIFAR-RT + MNIST-NRT & $88.10 \pm 0.06$ & $33.56 \pm 0.90$ \\
CIFAR-NRT + MNIST-RT & $91.26 \pm 0.05$ & $37.66 \pm 0.66$ \\
\bottomrule
\end{tabular}
\caption{MoE compositional robustness (averaged across router training methods). Expert configuration determines system robustness; router training shows $<$0.1\% difference across all metrics.}
\label{tab:moe_robustness}
\end{table}

\begin{table}[h]
\centering
\small
\begin{tabular}{lccc}
\toprule
\textbf{Router Training} & \textbf{Clean Acc} & \textbf{Adv Acc} & \textbf{Time (s)} \\
\midrule
RT & $89.64\%$ & $35.49\%$ & $4107$ \\
NRT & $89.70\%$ & $35.47\%$ & $968$ \\
\midrule
$\Delta$ & $+0.06\%$ & $-0.02\%$ & $\mathbf{4.2\times}$ faster \\
\bottomrule
\end{tabular}
\caption{Router training comparison (averaged across expert configurations). Adversarial router training provides no robustness benefit for dissimilar domains but incurs 4$\times$ training cost.}
\label{tab:router_training}
\end{table}

\textbf{Key Findings:}
\begin{enumerate}
    \item \textbf{Expert Dominance:} Expert configuration causes 13.15\% adversarial accuracy difference (Both RT: 41.93\% vs Both NRT: 28.78\%). Router training causes $<$0.1\% difference.
    
    \item \textbf{Perfect Routing:} 100\% gating accuracy across all 40 runs validates the compositional property---the router acts as a deterministic switch, faithfully propagating expert robustness.
    
    \item \textbf{Training Efficiency:} NRT routers achieve identical performance with 4$\times$ speedup for visually dissimilar domains.
    
    \item \textbf{Modular Tradeoffs:} Mixed configurations (e.g., CIFAR-NRT + MNIST-RT: 91.26\% clean, 37.66\% adv) enable domain-specific robustness requirements.
\end{enumerate}

\subsection{Empirical-Formal Correlation}

\begin{table}[h]
\centering
\small
\begin{tabular}{llccc}
\toprule
\textbf{Model} & \textbf{Training} & \textbf{AA (\%)} & \textbf{CRA (\%)} & $\boldsymbol{\Delta}$ \\
\midrule
CIFAR-10 & NRT & 0.01 & 0 & --- \\
CIFAR-10 & RT & 17.24 & 90 & 72.8\% \\
MNIST & NRT & 46.06 & 75 & 28.9\% \\
MNIST & RT & 90.77 & 100 & 9.2\% \\
Router & NRT/RT & 35.5 & 100 & 64.5\% \\
\bottomrule
\end{tabular}
\caption{Empirical-formal gap analysis at $\varepsilon = 2/255$. AA = Adversarial Accuracy (PGD), CRA = Certified Robustness Accuracy ($\alpha$--$\beta$-CROWN), $\Delta = \text{CRA} - \text{AA}$.}
\label{tab:gap_analysis}
\end{table}

\textbf{Interpretation:} CRA $\geq$ AA reflects formal verification providing worst-case guarantees while PGD provides heuristic lower bounds. The gap varies by domain: minimal for MNIST-RT (9.2\%), substantial for CIFAR-10-RT (72.8\%). Robust training consistently improves both metrics. The router's large gap (64.5\%) reflects geometric dataset separation creating decision margins $\gg \varepsilon$.

\subsection{Scalability}

\paragraph{Verification Efficiency.} Compositional verification achieves:
\begin{itemize}
    \item \textbf{Parameter reduction:} 2M (full MoE) $\to$ 96K (largest component) = 20$\times$
    \item \textbf{Verification time:} $\sim$4.6s per sample at $\varepsilon \leq 4/255$
    \item \textbf{Incremental integration:} New experts verified independently without full-system re-verification
\end{itemize}

\paragraph{Training Efficiency.} For 3-dataset MoE, compositional training (frozen experts + router fine-tuning) achieves 202 min vs 242.5 min for monolithic retraining---16.7\% reduction with greater savings as experts increase.

\paragraph{Inference Scaling.} MoE inference scales as $O(k)$ where $k$ is the number of activated experts, enabling efficient deployment: $k=1$ activation requires only router + single expert forward pass.

%==============================================================================
% 5. DISCUSSION AND LIMITATIONS
%==============================================================================
\section{Discussion and Limitations}

\paragraph{When Compositional Verification Applies.} Our framework requires router stability within the perturbation region. This holds for visually dissimilar domains (CIFAR-10 vs MNIST: 100\% routing accuracy) but degrades for similar domains. Experiments with traffic sign datasets (GTSRB + PTSD) show 93.17\% routing accuracy and 12\% adversarial accuracy reduction compared to dissimilar domains (Appendix~\ref{app:similar_domains}). Similar domains may require probabilistic certification or stronger router training.

\paragraph{Architecture Scope.} Our verification-compatible CNNs (96K parameters) enable tractable formal analysis but limit expressiveness. Scaling to transformers faces memory constraints with current verifiers. Promising directions include randomized smoothing for probabilistic certificates and verification-aware architecture search.

\paragraph{Practical Implications.} For safety-critical MoE deployment: (1) use adversarial training for experts (essential for robustness); (2) evaluate router training necessity based on domain dissimilarity; (3) compose verified components for system guarantees; (4) leverage modular structure for incremental expert integration.

%==============================================================================
% 6. CONCLUSION
%==============================================================================
\section{Conclusion}

We establish that robustness in heterogeneous MoE systems composes modularly: verified router stability combined with verified expert robustness yields verified system robustness. This decomposition reduces verification complexity by 20$\times$ and enables scalable modular certification.

Our empirical findings reveal clear design principles: expert robustness dominates system behavior (13.1\% impact) while router training shows negligible effect for dissimilar domains ($<$0.1\%), enabling significant training efficiency gains. Robust training is prerequisite for formal verifiability---non-robust experts are fundamentally incompatible with sound verification.

These results demonstrate that empirical robustness and formal verification form complementary facets of a unified framework, advancing the path toward trustworthy, compositional AI for safety-critical deployment.

%==============================================================================
% REFERENCES
%==============================================================================
\bibliographystyle{plain}
\begin{thebibliography}{10}

\bibitem{jacobs1991}
R.~A. Jacobs, M.~I. Jordan, S.~J. Nowlan, and G.~E. Hinton.
\newblock Adaptive mixtures of local experts.
\newblock {\em Neural Computation}, 3(1):79--87, 1991.

\bibitem{shazeer2017}
N.~Shazeer et~al.
\newblock Outrageously large neural networks: The sparsely-gated mixture-of-experts layer.
\newblock {\em ICLR}, 2017.

\bibitem{riquelme2021}
C.~Riquelme et~al.
\newblock Scaling vision with sparse mixture of experts.
\newblock {\em NeurIPS}, 2021.

\bibitem{fedus2023}
W.~Fedus, B.~Zoph, and N.~Shazeer.
\newblock Switch transformers: Scaling to trillion parameter models.
\newblock {\em JMLR}, 23(120):1--39, 2022.

\bibitem{pfeiffer2023}
J.~Pfeiffer et~al.
\newblock Modular deep learning.
\newblock {\em arXiv:2302.11529}, 2023.

\bibitem{madry2018}
A.~Madry et~al.
\newblock Towards deep learning models resistant to adversarial attacks.
\newblock {\em ICLR}, 2018.

\bibitem{wang2021}
S.~Wang et~al.
\newblock Beta-CROWN: Efficient bound propagation with per-neuron split constraints.
\newblock {\em NeurIPS}, 2021.

\end{thebibliography}

%==============================================================================
% APPENDIX
%==============================================================================
\newpage
\appendix

\section{Training Hyperparameters}
\label{app:hyperparams}

\textbf{Experts:} Adam optimizer, LR $10^{-3}$, 200 epochs, batch size 128, label smoothing 0.1. RT: PGD with $\varepsilon = 8/255$, 7 steps, step size $2/255$.

\textbf{Router:} Adam optimizer, LR $10^{-3}$, weight decay $5 \times 10^{-4}$, 50 epochs on frozen experts. Same PGD parameters for RT.

\section{Architecture Details}
\label{app:architecture}

\textbf{VerifiableCNN (96K params):}
\begin{itemize}
    \item Conv blocks: $3 \to 20 \to 28 \to 40 \to 56$ channels, $3 \times 3$ kernels, ReLU, AvgPool
    \item Global average pooling $\to$ FC(56, 128) $\to$ ReLU $\to$ FC(128, $C$)
    \item No BatchNorm (folded into conv weights for ONNX export)
\end{itemize}

\section{Impact of Dataset Similarity}
\label{app:similar_domains}

Experiments with GTSRB + PTSD (both traffic sign datasets) show routing accuracy degrades to 93.17\% (vs 100\% for CIFAR-10 + MNIST), and adversarial accuracy drops to 29.47\% (vs 41.59\%). The 6.83\% routing error propagates to 12.12\% system accuracy reduction, validating that domain dissimilarity enables the clean compositional guarantees in the main paper.

\section{Hardware and Software}
\label{app:hardware}

Intel i7-14700K, RTX 4060 8GB, 64GB RAM. PyTorch 2.1, $\alpha$--$\beta$-CROWN (Dec 2024), Gurobi 11.0. Verification timeout: 300s per sample.

\end{document}